% ============================================================
%  Глава 5. Эксперименты и результаты — sdf_raster
%  Черновик. Все [PLACEHOLDER] заменить реальными данными.
%  \cite{} и \ref{} — стандартные LaTeX-ключи.
% ============================================================

\chapter{Эксперименты и результаты}
\label{chap:experiments}

В данной главе представлены результаты экспериментальной оценки предложенного
метода растеризации SDF-сцен (\textbf{sdf\_raster}) на основе конвейера:
SCom $\to$ Frustum Culling $\to$ Hi-Z Occlusion Culling $\to$ Marching Cubes
(GPU compute) $\to$ Mesh Shaders $\to$ G-buffer $\to$ Deferred Shading.
Оцениваются производительность, качество геометрии, эффективность отсечения
и вклад отдельных компонентов конвейера.

% ------------------------------------------------------------
\section{Методология}
\label{sec:methodology}

% - - - - - - - - - - - - - - - - - - - - - - - - - - - - - -
\subsection{Аппаратная конфигурация}

Все эксперименты проводились на единой аппаратной платформе:
видеокарта~\textbf{[NVIDIA/AMD XXX]}, центральный процессор~\textbf{[XXX]}
с тактовой частотой~\textbf{[X.X]~ГГц}, оперативная память~\textbf{[XX]~ГБ
DDR[5]}, версия Vulkan~\textbf{[1.3.XXX]}, версия драйвера~\textbf{[XXX.XX]}.
Использование единой аппаратной конфигурации исключает аппаратные факторы
как источник вариативности результатов.

% - - - - - - - - - - - - - - - - - - - - - - - - - - - - - -
\subsection{Тестовые сцены}

В качестве тестовых объектов использовались четыре классические модели
из Stanford 3D Scanning Repository~\cite{stanford_models}
(Таблица~\ref{tab:scenes}).

\begin{table}[h]
  \centering
  \caption{Тестовые сцены и диапазон размеров файлов SCom.}
  \label{tab:scenes}
  \begin{tabular}{lrr}
    \toprule
    \textbf{Модель} & \textbf{Треугольников (mesh)} & \textbf{Размер SCom, МБ} \\
    \midrule
    Dragon       &   871\,000 & [X.X]--[XX] \\
    Bunny        &    69\,000 & [X.X]--[XX] \\
    Armadillo    &   346\,000 & [X.X]--[XX] \\
    Happy Buddha & 1\,087\,000 & [X.X]--[XX] \\
    \bottomrule
  \end{tabular}
\end{table}

Каждая модель кодировалась в формат SCom~\cite{scom2025} при нескольких
уровнях сжатия, что позволяет варьировать размер файла и систематически
исследовать зависимость производительности и качества от степени компрессии.
Диапазон размеров охватывает более двух порядков величины~---
от~\textbf{[X]~МБ} до~\textbf{[XX]~МБ}.
Объём внутренних структур в видеопамяти в общем случае превышает размер файла
в~\textbf{[X.X]--[X.X]~раза}.

% - - - - - - - - - - - - - - - - - - - - - - - - - - - - - -
\subsection{Метрики}

\paragraph{Частота кадров (FPS).}
Среднее значение FPS вычислялось по~1\,000 последовательным кадрам после
прогрева (100~кадров), исключающего переходные процессы инициализации кэша GPU.
Временны́е метки фиксировались GPU timestamp queries интерфейса Vulkan
(\texttt{VK\_QUERY\_TYPE\_TIMESTAMP}), что гарантирует измерение реального
времени GPU без вмешательства CPU.

\paragraph{Пиковое отношение сигнал/шум (PSNR).}
PSNR оценивалось относительно эталонного изображения, сформированного
в \emph{собственном рендерере} при плотной тесселяции
(${\approx}10\,000\,000$~треугольников, тот же G-buffer + Blinn–Phong конвейер).
Такой выбор изолирует ошибку Marching Cubes от различий в модели освещения,
делая PSNR чистым показателем геометрической точности:
\[
  \mathrm{PSNR} = 10 \cdot \log_{10}\!\left(\frac{255^2}{\mathrm{MSE}}\right),
\]
где MSE рассчитывалось по всем пикселям кадра в пространстве RGB.

\paragraph{Эффективность отсечения.}
Доля узлов октодерева, исключённых Frustum Culling и Hi-Z Occlusion Culling,
от общего числа~--- атомарные счётчики в шейдерах и
pipeline statistics queries.

\paragraph{Разбивка времени кадра.}
Независимое время GPU для каждого из семи проходов:
(1)~Frustum Culling compute,
(2)~Hi-Z generation,
(3)~Hi-Z Occlusion Culling compute,
(4)~Marching Cubes compute,
(5)~Mesh Shader (task\,+\,mesh),
(6)~G-buffer pass,
(7)~Deferred Shading.

% - - - - - - - - - - - - - - - - - - - - - - - - - - - - - -
\subsection{Базовые методы сравнения}

\begin{table}[h]
  \centering
  \caption{Базовые методы сравнения.}
  \label{tab:baselines}
  \begin{tabular}{lp{6cm}p{4.5cm}}
    \toprule
    \textbf{Метод} & \textbf{Описание} & \textbf{Цель сравнения} \\
    \midrule
    SCom Ray Tracing
      & Исходная реализация SCom~\cite{scom2025} с аппаратной трассировкой
        лучей (\texttt{VK\_KHR\_ray\_tracing\_pipeline})
      & Прямое сравнение RT vs.\ растеризация \\[4pt]
    Naive MC + Vertex Pipeline
      & Полный MC на GPU, передача меша через \texttt{vkCmdDraw}
        без Mesh Shaders и иерархического отсечения
      & Нижняя оценка; вклад Mesh Shaders \\[4pt]
    Frustum-only MC
      & Frustum Culling активирован, Hi-Z Occlusion Culling отключён
      & Изолированный вклад Hi-Z \\[4pt]
    \textbf{sdf\_raster (ours)}
      & Полный конвейер: Frustum\,+\,Hi-Z\,+\,MC\,+\,Mesh Shaders
      & Предложенный метод \\
    \bottomrule
  \end{tabular}
\end{table}

% ============================================================
\section{Производительность}
\label{sec:performance}

% - - - - - - - - - - - - - - - - - - - - - - - - - - - - - -
\subsection{Основные результаты}

Таблица~\ref{tab:main_results} представляет сводные результаты
производительности (FPS) для всех тестовых сцен и базовых методов.
Рисунок~\ref{fig:fps_vs_size} иллюстрирует зависимость частоты кадров от
размера модели SCom (МБ) в логарифмическом масштабе по оси абсцисс.

\begin{table}[h]
  \centering
  \caption{Производительность (FPS) для четырёх тестовых сцен при трёх
    уровнях сжатия SCom и четырёх методах рендеринга.}
  \label{tab:main_results}
  \begin{tabular}{lrrrr}
    \toprule
    \textbf{Сцена / МБ}
      & \textbf{SCom RT}
      & \textbf{Naive MC}
      & \textbf{Frustum-only}
      & \textbf{sdf\_raster} \\
    \midrule
    Dragon / [X] МБ       & [XX.X] & [XX.X] & [XX.X] & [XXX]   \\
    Dragon / [XX] МБ      & [XX.X] & [XX.X] & [XX.X] & [XX.X]  \\
    Bunny / [X] МБ        & [XX.X] & [XX.X] & [XX.X] & [XXX]   \\
    Bunny / [XX] МБ       & [XX.X] & [XX.X] & [XX.X] & [XX.X]  \\
    Armadillo / [X] МБ    & [XX.X] & [XX.X] & [XX.X] & [XXX]   \\
    Armadillo / [XX] МБ   & [XX.X] & [XX.X] & [XX.X] & [XX.X]  \\
    Buddha / [X] МБ       & [XX.X] & [XX.X] & [XX.X] & [XXX]   \\
    Buddha / [XX] МБ      & [XX.X] & [XX.X] & [XX.X] & [XX.X]  \\
    \bottomrule
  \end{tabular}
\end{table}

\begin{figure}[h]
  \centering
  % \includegraphics[width=\linewidth]{figures/fps_vs_size.pdf}
  \fbox{\parbox{0.9\linewidth}{\centering\vspace{2cm}
    [РИСУНОК: FPS vs.\ Размер модели SCom (МБ, лог.\ шкала).
     Кривые: sdf\_raster, SCom RT, Naive MC, Frustum-only.]
  \vspace{2cm}}}
  \caption{FPS vs.\ Размер модели SCom (МБ, логарифмическая шкала).}
  \label{fig:fps_vs_size}
\end{figure}

Предложенный метод \textbf{sdf\_raster} обеспечивает
\textbf{[XX.X]--[XXX]~FPS} в диапазоне от наиболее крупной до наиболее
компактной модели. SCom Ray Tracing демонстрирует \textbf{[XX.X]--[XXX]~FPS},
что на \textbf{[X]--[X] раз} ниже на крупных моделях и сопоставимо на малых.
Таким образом, растеризационный конвейер превосходит трассировщик лучей
в среднем в~\textbf{[X.X]~раза} при сопоставимом качестве изображения
(см.\ раздел~\ref{sec:quality}).

Наивный конвейер (Naive MC\,+\,Vertex Pipeline) показывает
\textbf{[XX.X]--[XX.X]~FPS}. Ключевым узким местом является
синхронизационный барьер CPU--GPU при передаче геометрии через
\texttt{vkCmdDraw}. Применение Mesh Shaders устраняет данное узкое место
и обеспечивает прирост~\textbf{[XX]\%}.

Frustum-only MC занимает промежуточное положение:
\textbf{[XX.X]--[XX.X]~FPS}.
Добавление Hi-Z Occlusion Culling даёт дополнительный прирост
\textbf{[XX]\%}, наиболее выраженный в сценах с высоким взаимным
перекрытием геометрии (Dragon, Armadillo).

% - - - - - - - - - - - - - - - - - - - - - - - - - - - - - -
\subsection{Масштабируемость и точка пересечения}

Логарифмическая шкала по оси абсцисс выбрана намеренно: диапазон размеров
охватывает более двух порядков величины, и линейная шкала привела бы к
скоплению точек. Характерной особенностью является \textbf{точка пересечения}
кривых sdf\_raster и Naive MC при размере модели ${\approx}$\textbf{[X.X]~МБ}:
при меньших размерах накладные расходы на Mesh Shaders и иерархическое
отсечение нивелируют их преимущества. При больших размерах преимущество
sdf\_raster возрастает монотонно, согласуясь с теоретической
субквадратичной сложностью иерархического отсечения~\cite{greene1993hier}.

% - - - - - - - - - - - - - - - - - - - - - - - - - - - - - -
\subsection{Разбивка времени кадра}

\begin{table}[h]
  \centering
  \caption{Разбивка времени GPU-кадра по этапам конвейера
    (сцена [SCENE\_NAME], [XX]~МБ).}
  \label{tab:frame_breakdown}
  \begin{tabular}{lrr}
    \toprule
    \textbf{Этап конвейера} & \textbf{Время, мс} & \textbf{Доля, \%} \\
    \midrule
    Frustum Culling compute  & [X.XX] & [XX] \\
    Hi-Z generation          & [X.XX] & [XX] \\
    Hi-Z Occlusion Culling   & [X.XX] & [XX] \\
    Marching Cubes compute   & [X.XX] & [XX] \\
    Mesh Shader (task+mesh)  & [X.XX] & [XX] \\
    G-buffer pass            & [X.XX] & [XX] \\
    Deferred Shading         & [X.XX] & [XX] \\
    \midrule
    \textbf{Итого}           & \textbf{[X.XX]} & \textbf{100} \\
    \bottomrule
  \end{tabular}
\end{table}

Доминирующим этапом является Marching Cubes compute (\textbf{[X.XX]~мс})
из-за нерегулярного доступа к SDF-кирпичам. Проход Mesh Shader~---
второй по значимости (\textbf{[X.XX]~мс}). Frustum Culling и Hi-Z Occlusion
Culling совместно занимают менее~\textbf{[XX]\%} общего времени кадра, при
этом обеспечивая значительное сокращение работы на последующих этапах.

% ============================================================
\section{Качество геометрии}
\label{sec:quality}

% - - - - - - - - - - - - - - - - - - - - - - - - - - - - - -
\subsection{PSNR и зависимость от уровня сжатия}

Рисунок~\ref{fig:psnr_vs_size} показывает зависимость PSNR~(дБ) от размера
модели SCom~(МБ). Увеличение степени сжатия SCom огрубляет октодерево и
снижает число кирпичей, кодирующих тонкие детали. При высоком сжатии
(\textbf{[X.X]~МБ}) PSNR составляет~\textbf{[XX.X]~дБ}~--- визуально
заметное сглаживание. При низком сжатии (\textbf{[XX]~МБ}) PSNR достигает
\textbf{[XX.X]~дБ}, неотличимого от эталона при стандартном игровом
расстоянии.

\begin{figure}[h]
  \centering
  % \includegraphics[width=\linewidth]{figures/psnr_vs_size.pdf}
  \fbox{\parbox{0.9\linewidth}{\centering\vspace{2cm}
    [РИСУНОК: PSNR~(дБ) vs.\ Размер модели SCom~(МБ).
     Кривые: sdf\_raster (ours) и SCom Ray Tracing.]
  \vspace{2cm}}}
  \caption{PSNR~(дБ) vs.\ Размер модели SCom~(МБ).}
  \label{fig:psnr_vs_size}
\end{figure}

\textbf{Важно:} в качестве эталона используется \emph{не} сторонний
офлайн-рендерер (Blender Cycles, PBRT), а тот же конвейер G-buffer\,+\,Deferred
Shading собственного рендерера, применённый к плотно тесселированному мешу
(${\approx}10\,000\,000$~треугольников). Это решение изолирует ошибку
Marching Cubes от различий в модели освещения~--- PSNR становится
\textbf{чистым показателем геометрической точности}.

% - - - - - - - - - - - - - - - - - - - - - - - - - - - - - -
\subsection{Сравнение с SCom Ray Tracing}

SCom Ray Tracing~\cite{scom2025} получает PSNR~\textbf{[XX.X]--[XX.X]~дБ}.
Трассировщик лучей точнее воспроизводит изоповерхность SDF: каждый луч
находит пересечение аналитически в пространстве кирпича без дискретизации.
Растеризационный путь sdf\_raster уступает в среднем на~\textbf{[X.X]~дБ}
(известное ограничение Marching Cubes~--- ошибка дискретизации, пропорциональная
шагу сетки). Тем не менее разница \textbf{[X.X]~дБ} практически незаметна
на репрезентативных изображениях (Рис.~\ref{fig:visual_comparison})
за исключением максимально приближённых ракурсов у острых рёбер.

% - - - - - - - - - - - - - - - - - - - - - - - - - - - - - -
\subsection{Визуальные артефакты}

\begin{figure}[h]
  \centering
  % \includegraphics[width=\linewidth]{figures/visual_comparison.png}
  \fbox{\parbox{0.9\linewidth}{\centering\vspace{2cm}
    [РИСУНОК: Визуальное сравнение~--- sdf\_raster (слева),
     SCom Ray Tracing (центр), эталон --- плотный меш (справа).
     Сцена Dragon, ракурс у острого ребра, 1920$\times$1080.]
  \vspace{2cm}}}
  \caption{Визуальное сравнение: sdf\_raster (слева), SCom Ray Tracing
    (центр), эталон~--- плотный меш (справа).}
  \label{fig:visual_comparison}
\end{figure}

Наиболее выраженный артефакт~--- \textbf{ступенчатость вдоль острых рёбер}
при низких уровнях сжатия. Вычисление нормалей по градиенту SDF в вершинах
Marching Cubes смягчает артефакт в закраске, не улучшая геометрию.
При высоких уровнях сжатия (большие файлы МБ) артефакт исчезает.

% ============================================================
\section{Эффективность отсечения}
\label{sec:culling}

% - - - - - - - - - - - - - - - - - - - - - - - - - - - - - -
\subsection{Результаты по сценам}

\begin{table}[h]
  \centering
  \caption{Эффективность отсечения: доля отсечённых узлов октодерева~(\%).}
  \label{tab:culling_efficiency}
  \begin{tabular}{lrrr}
    \toprule
    \textbf{Сцена}
      & \textbf{Frustum Culling, \%}
      & \textbf{Hi-Z, \%}
      & \textbf{Суммарно, \%} \\
    \midrule
    Dragon       & [XX] & [XX] & [XX] \\
    Bunny        & [XX] & [XX] & [XX] \\
    Armadillo    & [XX] & [XX] & [XX] \\
    Happy Buddha & [XX] & [XX] & [XX] \\
    \bottomrule
  \end{tabular}
\end{table}

Для Frustum Culling доля отсекаемых узлов составляет~\textbf{[XX]--[XX]\%}
и зависит от соотношения объёма ограничивающего AABB к объёму пирамиды
видимости (FOV\,=\,\textbf{[XX]}°). Модели с высоким аспектным отношением
(Dragon, Happy Buddha) демонстрируют более высокий процент отсечения.

Hi-Z Occlusion Culling добавляет \textbf{[XX]--[XX]\%} к уже отсечённым
узлам, обеспечивая суммарное отсечение~\textbf{[XX]--[XX]\%}.
Максимальный выигрыш Hi-Z наблюдается в закрытых ракурсах (Dragon,
Armadillo), где передняя геометрия перекрывает значительную долю задней.

% - - - - - - - - - - - - - - - - - - - - - - - - - - - - - -
\subsection{Накладные расходы на отсечение}

Генерация Hi-Z буфера занимает~\textbf{[X.XX]~мс}; построение иерархической
пирамиды глубины (Hi-Z mip chain)~--- \textbf{[X.XX]~мс}, практически не
зависит от сложности сцены. Отношение затрат на отсечение к сэкономленным
вычислениям свидетельствует о высокой рентабельности для сцен с числом
узлов октодерева более~\textbf{[XXX]}.

% ============================================================
\section{Анализ вклада компонентов (Ablation Study)}
\label{sec:ablation}

\begin{table}[h]
  \centering
  \caption{Ablation study: FPS при пошаговом добавлении компонентов конвейера
    (сцена [SCENE\_NAME], [XX]~МБ).}
  \label{tab:ablation}
  \begin{tabular}{lrr}
    \toprule
    \textbf{Конфигурация} & \textbf{FPS} & \textbf{Прирост, \%} \\
    \midrule
    Naive MC + Vertex Pipeline (baseline) & [XX.X] & ---    \\
    + Frustum Culling                     & [XX.X] & +[XX]  \\
    + Hi-Z Occlusion Culling              & [XX.X] & +[XX]  \\
    + Mesh Shaders (task\,+\,mesh)        & [XX.X] & +[XX]  \\
    \textbf{sdf\_raster (full pipeline)}  & \textbf{[XX.X]} & \textbf{+[XX]} \\
    \bottomrule
  \end{tabular}
\end{table}

Базовая конфигурация (Naive MC\,+\,Vertex Pipeline) обеспечивает
\textbf{[XX.X]~FPS}. Добавление Frustum Culling повышает FPS до
\textbf{[XX.X]} (\textbf{+[XX]\%}). Включение Hi-Z Occlusion Culling~---
до~\textbf{[XX.X]} (\textbf{+[XX]\%}). Замена вершинного конвейера на Mesh
Shaders обеспечивает наибольший единичный прирост:
\textbf{[XX.X]~FPS (+[XX]\%)}~--- устранение CPU-узкого места является
доминирующим фактором масштабируемости. Полный конвейер достигает
\textbf{[XX.X]~FPS}~--- в~\textbf{[X.X]~раза} выше наивной реализации.

\textbf{Практическое правило:} на моделях малого размера ([X]--[X]~МБ)
вклад Hi-Z пренебрежимо мал (менее~[X]\%), поскольку число узлов октодерева
невелико и накладные расходы на Hi-Z-пирамиду не окупаются. При размере
модели [XX]~МБ и выше преимущество Hi-Z становится значимым и монотонно
возрастает.

% ============================================================
\section{Ограничения}
\label{sec:limitations}

\paragraph{Артефакты дискретизации Marching Cubes.}
При низких уровнях сжатия алгоритм вводит заметную ступенчатость вдоль
острых рёбер. Ни постобработка нормалей, ни увеличение разрешения сетки MC
не устраняют артефакт полностью. Альтернатива~--- адаптивные схемы, например
Dual Contouring~\cite{ju2002dualcontouring},~--- нетривиальна в интеграции с
Mesh Shaders.

\paragraph{Деградация производительности при экстремальных размерах.}
При размере файла~\textbf{[XX]~МБ} метод обеспечивает~\textbf{[X]~FPS}~---
приемлемо для интерактивного просмотра, но недостаточно для видеоигр
(целевые 60~FPS). Число запусков Marching Cubes достигает предельных значений
доступной видеопамяти.

\paragraph{Отсутствие поддержки динамических SDF-сцен.}
Текущая реализация предполагает статичность SCom-модели: изменение SDF
требует полного перекодирования на CPU и повторной выгрузки в VRAM.
Поддержка анимированных объектов является открытой проблемой.

\paragraph{Потребление видеопамяти.}
Развёртывание SCom-модели~\textbf{[XX]~МБ} требует~\textbf{[XXX]~МБ VRAM}.
Доминирующий вклад~--- разжатые SDF-кирпичи (16-разрядный FP). Приемлемо
для GPU с VRAM от~\textbf{[X]~ГБ}, но может быть ограничивающим фактором на
мобильных/встроенных платформах.

% ============================================================
\section*{Краткое резюме главы}

В данной главе представлены результаты комплексной экспериментальной оценки
метода \textbf{sdf\_raster}. Ключевые выводы:

\begin{itemize}
  \item \textbf{Производительность:} sdf\_raster превосходит SCom Ray
    Tracing~\cite{scom2025} в среднем в~\textbf{[X.X]~раза} на крупных
    моделях при незначительном снижении PSNR на~\textbf{[X.X]~дБ}.
  \item \textbf{Mesh Shaders:} замена вершинного конвейера на Mesh Shaders
    обеспечивает наибольший единичный прирост производительности, устраняя
    CPU-узкое место.
  \item \textbf{Hi-Z Occlusion Culling:} иерархическое Hi-Z добавляет
    \textbf{[XX]\%} прироста на сценах с высоким перекрытием геометрии при
    минимальных накладных расходах.
  \item \textbf{Качество:} PSNR достигает \textbf{[XX.X]~дБ} при высоких
    уровнях детализации SCom, визуально неотличимых от эталона на стандартном
    расстоянии просмотра.
  \item \textbf{Ablation Study:} все компоненты конвейера вносят измеримый
    вклад; их совокупный эффект превышает сумму вкладов отдельных компонентов.
\end{itemize}

Выявленные ограничения~--- прежде всего артефакты Marching Cubes и
отсутствие поддержки анимации~--- намечают направления дальнейших исследований.

% ============================================================
%  СПРАВОЧНИК ПЛЕЙСХОЛДЕРОВ
% ============================================================
%
%  [NVIDIA/AMD XXX]   — модель GPU
%  [XXX] (CPU)        — модель CPU
%  [X.X] ГГц         — тактовая частота CPU
%  [XX] ГБ           — объём RAM
%  [1.3.XXX]          — версия Vulkan
%  [XXX.XX]           — версия драйвера
%  [X]–[XX] МБ       — диапазон размеров SCom-файлов
%  [X.X]–[X.X] раза  — коэффициент расширения в VRAM
%  [XX.X] FPS         — FPS для конкретного метода/сцены
%  [XXX] FPS          — максимальный FPS
%  [X.X] раза         — кратный прирост производительности
%  [XX]%              — процентный прирост FPS
%  [X.XX] мс          — время этапа GPU
%  [SCENE_NAME]       — имя тестовой сцены
%  [XX.X] дБ          — PSNR
%  [X.X] дБ           — разница PSNR между методами
%  [XXX] (узлы)       — пороговое число узлов октодерева
%  [XXX] МБ (VRAM)    — потребление видеопамяти
%  [X] ГБ (VRAM)      — минимальный объём VRAM
%
% ============================================================
%  BIBTEX-КЛЮЧИ
% ============================================================
%
%  \cite{scom2025}            — SCom: Structured SDF Compression, SIGGRAPH Asia 2025
%  \cite{stanford_models}     — Stanford 3D Scanning Repository
%  \cite{greene1993hier}      — Greene et al., Hierarchical Z-Buffer Visibility, SIGGRAPH 1993
%  \cite{ju2002dualcontouring} — Ju et al., Dual Contouring of Hermite Data, SIGGRAPH 2002
%
